\documentclass[12pt]{article}
\usepackage[utf8]{inputenc}
\usepackage[T1]{fontenc}
\usepackage[margin=1in]{geometry}
\usepackage{setspace}
\usepackage{xcolor}
\usepackage[colorlinks=true, urlcolor=blue!60!black, linkcolor=blue!60!black]{hyperref}
\setlength{\parindent}{2em}
\setlength{\parskip}{1.25em}
\renewcommand{\baselinestretch}{1.0}
\begin{document}
\begin{center}
{\LARGE \textit{Nixon v. United States}, 506 U.S. 224 (1993)\par}
\vspace{0.8em}
{\large Case Brief --- Four-Page Edition\par}
\end{center}
\vspace{1em}

\section{Memory Jogger}

Because Article I, \S~3, cl.~6 gives the Senate the ``sole Power to try all Impeachments'' and the word ``try'' supplies no judicially manageable standard, a removed judge's challenge to the Senate's use of a factfinding committee is a nonjusticiable political question (\href{https://www.courtlistener.com/opinion/112803/nixon-v-united-states/}{CourtListener}).

\section{Detailed Case Facts}

Walter L. Nixon, Jr., Chief Judge of the Southern District of Mississippi, was convicted of two counts of making false statements to a federal grand jury investigating whether he had accepted a gratuity to ask a district attorney to drop a prosecution of a businessman's son. He refused to resign and kept drawing his salary in prison. On May 10, 1989, the House adopted three articles of impeachment.

The Senate invoked Impeachment Rule XI, which lets a committee of Senators ``receive evidence and take testimony.'' The committee held four days of hearings with ten witnesses, including Nixon, and gave the full Senate a transcript and report. After final briefs and floor argument, in which Nixon made a personal appeal and Senators questioned both sides, the Senate convicted him on two articles by more than two-thirds on November 3, 1989, and removed him.

\section{Procedural History}

Nixon sued in the District of Columbia for a declaration that his conviction was void and for reinstatement of his salary, arguing that Rule XI prevented the full Senate from ``trying'' him. The district court dismissed the case as a nonjusticiable political question, and the D.C. Circuit affirmed. The Supreme Court granted certiorari and affirmed on January 13, 1993.

\section{Judicial Votes}

Affirmed, \textbf{9--0}, in four opinions.
\begin{itemize}
\item \textbf{Majority:} Chief Justice Rehnquist, joined by Justices Stevens, O'Connor, Scalia, Kennedy, and Thomas.
\item \textbf{Concurrence:} Justice Stevens.
\item \textbf{Concurring in the judgment:} Justice White, joined by Justice Blackmun (justiciable, but Rule XI is valid); Justice Souter (prudential nonjusticiability).
\item \textbf{Dissents:} None.
\end{itemize}

\section{Holding}

A challenge to the Senate's use of a committee to take evidence under Rule XI is nonjusticiable: the Impeachment Trial Clause textually commits the issue to the Senate, and ``try'' lacks judicially manageable standards. The Court's broader remarks about presidential impeachments, finality, and remedies go beyond the facts. By distinguishing \textit{Powell v. McCormack} on the basis of ``identifiable textual limits,'' the Court implied that violating the oath, two-thirds, or Chief Justice requirements might remain reviewable.

\section{Analysis of Opinions}

\textbf{Majority (Rehnquist).} Applying the first two \textit{Baker v. Carr} factors, which reinforce each other, Rehnquist found that dictionaries from Johnson's to Webster's give ``try'' broad meanings, and that the Clause's three specific requirements show the Framers omitted others. ``Sole,'' used only in the two Impeachment Clauses, means functioning independently and without interference; the Convention adopted the final text, so the Committee of Style's addition counts. The Convention rejected trials in the Supreme Court, and \textit{The Federalist} described impeachment as the Legislature's check on the Judiciary, so judicial review would let judges police the check on themselves, threaten finality, and pose intractable remedial problems. \textit{Powell} differed because separate text fixed House qualifications.

\textbf{Stevens (concurring).} The Framers trusted the Legislature, and history since Justice Chase's acquittal refutes the ``improbable hypotheticals'' of White and Souter.

\textbf{White (concurring in the judgment).} The Solicitor General conceded the Senate could convict an official as simply ``a bad guy'' without any hearing. ``Sole'' divides power between the Houses, and ``try'' is as manageable as ``due process of law.'' On the merits, English practice, Jefferson's \textit{Manual}, Story, and special masters show committee factfinding is consistent with a trial.

\textbf{Souter (concurring in the judgment).} Prudential \textit{Baker} factors bar review here, but courts might intervene if the Senate convicted ``upon a coin toss.''

\section{Examples: Future Applications}

\begin{enumerate}
\item \textbf{Same-Side (Nonjusticiable): Trial Rules.} Fifteen-minute cross-examination limits and hearsay exclusions in a Senate trial are the Senate's to design.
\item \textbf{Same-Side (Nonjusticiable): House Procedures.} A party-line House impeachment without hearings is unreviewable, because the House has the parallel ``sole Power of Impeachment.''
\item \textbf{Opposite-Side (Justiciable): Exclusion.} Excluding a Representative who meets the age, citizenship, and residency requirements is reviewable under \textit{Powell}, which \textit{Nixon} preserved.
\item \textbf{Opposite-Side (Justiciable): Majority Conviction.} A 55--45 conviction violates the express requirement of ``the Concurrence of two thirds of the Members present'' and should be reviewable.
\item \textbf{Fence-Sitter: Summary Conviction.} A conviction on a five-minute speech with no hearing might be unreviewable under Rehnquist's reasoning, yet White's and Souter's approaches would allow intervention.
\end{enumerate}

\section{Critique}

\textbf{Textualist views.} Many praise the opinion for respecting the Constitution's allocation of authority and for insulating Congress's only removal check on judges from judicial self-protection. Others echo White: ``sole'' allocated power between the Houses, and declaring a constitutional verb unenforceable risks turning limits into parchment barriers.

\textbf{Doctrinal views.} Scholars note tension with \textit{Marbury}, leaving impeachment to ``customary'' norms enforced by politics rather than courts. Souter's carve-out is unstable, since deciding whether the Senate stayed within ``broad boundaries'' already requires a merits judgment. The case pits rule-of-law accountability against separation-of-powers finality.

\section{Key Quotations}

\begin{quote}
``The conclusion that the use of the word `try' in the first sentence of the Impeachment Trial Clause lacks sufficient precision to afford any judicially manageable standard of review of the Senate's actions is fortified by the existence of the three very specific requirements that the Constitution does impose on the Senate when trying impeachments: The Members must be under oath, a two-thirds vote is required to convict, and the Chief Justice presides when the President is tried.''
\end{quote}
\noindent --- Chief Justice Rehnquist, Opinion of the Court, 506 U.S. at 230.

\begin{quote}
``But we conclude, after exercising that delicate responsibility, that the word `try' in the Impeachment Trial Clause does not provide an identifiable textual limit on the authority which is committed to the Senate.''
\end{quote}
\noindent --- Chief Justice Rehnquist, Opinion of the Court, 506 U.S. at 238.

\begin{quote}
``Judicial involvement in impeachment proceedings, even if only for purposes of judicial review, is counterintuitive because it would eviscerate the `important constitutional check' placed on the Judiciary by the Framers.''
\end{quote}
\noindent --- Chief Justice Rehnquist, Opinion of the Court, 506 U.S. at 235.

\end{document}
